\scanpage{11}
santes $X_\alpha$ de genre \inserted{$g_\alpha$} donné (i.e. des $\widetilde X_\alpha$ de genre $g_\alpha\ldots$). De plus, quand on se donne \deleted{\unsure{des}} $I\to S$ ($I$ ens fini), cela peut être réalisé par un $I\subset\widehat X^{\mathrm{lisse}}$, i.e. par une courbe \unsure{standard|standart} $S$.

La \uline{maquette} d'une courbe \unsure{standard|standart} $X$ consiste, par définition, en les données suivantes
\[
(10)\quad\left\{\begin{aligned}
&\text{a) le graphe }G_X=\biggl(
\underset{\inserted{\substack{\Vert\\\pi_0(\widetilde X)\\\Vert\\\pi_0(\widehat{\widetilde X})}}}{S},
\underset{\inserted{\substack{\Vert\\\unsure{\operatorname{Sing}(X)}}}}{\widetilde A},
\underset{\deleted{\unsure{\widehat X}}}{\sigma_{\widetilde A}}
\biggr)\\
&\text{b) L'}\deleted{\text{application}}\inserted{\text{ensemble}}\ I=\widehat X\setminus X,\text{ et l'application}\\
&\qquad I\longrightarrow S\\
&\text{c) L'application}\\
&\qquad S\xrightarrow{\uline{g}}\mathbb N\\
&\qquad\alpha\longmapsto\text{genre de }Y_\alpha=\widehat{\widetilde X}_\alpha
\end{aligned}\right.
\]
Une structure formée d'un graphe fini $G$\inserted{$=(S,\widetilde A,\sigma_{\widetilde A})$}, d'un ensemble fini $I$ \deleted{\unsure{de}} au dessus de l'ens. des sommets $S$ de $G$, et d'une application \inserted{``genre''} : $S\xrightarrow{\uline{g}}\mathbb N$, sera
\mspage{10}
